\documentclass[a4paper]{article}
\usepackage[pdftex]{graphicx}
\usepackage[utf8]{inputenc}
\usepackage{enumerate}
\usepackage{icomma}
\usepackage{amssymb}
\usepackage{tikz}
\usepackage{href-ul}
\hypersetup{
	colorlinks=true,
	linkcolor=blue,
	urlcolor=blue}
\usepackage{geometry}
\geometry{a4paper, top=15mm, left=15mm, right=15mm, bottom=15mm,
	headsep=10mm, footskip=12mm}

\begin{document}
\begin{tikzpicture}(10,3)
	\draw[ultra thick](2,0) --(10,0) -- (10,1.3) --(2,1.3) -- (2,0);
	\draw[fill=black](2,0)-- (0,.6) -- (2,.6) -- (2,0);
	\node at (6,.6) {\large Lösungsblatt von \href{https://www.okuyakl.de}{www.okuyakl.de}};
\end{tikzpicture}
\vspace{0.5 cm}

\noindent{\bf Aufgabe 1.}\\
Kettenregel: 
$$f'(x)= {1 \over 2} (5x^2+1)^{-{1 \over 2}} \cdot 10x$$

\noindent {\bf Aufgabe 1. b)}\\
Produkt- und Kettenregel:
$$g'(x)=ln(9-x^2) - x \cdot {2x \over 9-x^2}$$

\noindent {\bf Aufgabe 2. a)}\\
Der Numerus des Logarithmus muss größer als Null sein:
$$2x-3 > 0 \quad \Rightarrow \quad x > 1,5 $$

$$ \mathbb{D}_f= [1,5; \infty [  $$
Diese stetige Funktion untersuchen wir auf Grenzwerte:
$$
\renewcommand{\arraystretch}{1.5}
\begin{array}{rcl}
\lim\limits_{x\to 1,5}{\ln(2x-3)}&=& -\infty \\
\lim\limits_{x\to \infty}{\ln(2x-3)}&=& +\infty \\ 
\Rightarrow \mathbb{W}_f &=& \mathbb{R}
\end{array}
$$

\noindent {\bf Aufgabe 2. b)}\\
Wir berechnen zuerst die Nullstelle des Graphen von f:
$$
\renewcommand{\arraystretch}{1.5}
\begin{array}{rcll}
\ln(2x-3)&=& 0 &| e^{(\dots)}\\
   2x-3  &=& e^0=1  &| +3 \\ 
    2x   &=& 4      &| :2 \\
    x_0  &=& 2
\end{array}
$$
Dann leiten wir die Funktion ab: 
$$ f'(x)={2 \over 2x-3}$$ 
Eine Tangente erfüllt die Geradengleichung 
$$ y= m\cdot x + t $$
mit  
$$
\renewcommand{\arraystretch}{1.5}
\begin{array}{rcccl}
       y &=& f(2) &=&  0 \\
       m &=& f'(2)&=& {2 \over 4-3}=2 \\ 
y=mx+t &\Rightarrow & 0=2\cdot 2 +t & & |-4 \\
 & & t &=& -4 
\end{array}
$$

Die Gleichung der Tangente ist $$y=2x-4$$

\noindent {\bf Aufgabe 2. c)}\\
Die Funktion ist in ihrem Definitionsbereich streng monoton steigend, weil gilt: 
$$f'(x)>0 \quad \forall \quad \textnormal{(für alle)} \quad x \in \mathbb{D}_f$$
Begründung: Ein Bruch nimmt den Wert Null an, wenn der Zähler null wird. Im Term 
$ f'(x)={2 \over 2x-3}$ kann der Zähler jedoch nicht null werden, egal, was man für x einsetzt.
Die Umkehrfunktion lässt sich wie folgt bestimmen:
$$
\renewcommand{\arraystretch}{1.5}
\begin{array}{rcll}
\ln(2x-3)&=& y      &| e^{(\dots)}\\
   2x-3  &=& e^y    &| +3 \\ 
    2x   &=& e^y +3 &| :2 \\
      x  &=& {e^y +3 \over 2} &| x\rightleftharpoons y\\
f^{-1}(x)&=& {e^x +3 \over 2}
\end{array}
$$
Die Definitionsmenge einer Umkehrfunktion ist gleich der Wertemenge der ursprünglichen Funktion,
also $$\mathbb{D}_{f^{-1}}= \mathbb{R}$$

\noindent {\bf Aufgabe 3. a)}\\
Wir leiten $f(x)$ mit der Produkt- und Kettenregel ab:
$$f'(x)=e^{1-x} + x \cdot e^{1-x} \cdot (-1) = (1-x) \cdot e^{1-x}$$

\noindent {\bf Aufgabe 3. b)}\\
Für die Untersuchung des Monotonieverhaltens setzen wir zunächst $f'(x)=0$:

$$
\renewcommand{\arraystretch}{1.5}
\begin{array}{rcll}
(1-x) \cdot e^{1-x} &=& 0   &| :e^{1-x} \quad e^{(\dots)} \quad \textnormal{kann nie null werden}\\
               1-x  &=& 0   &| +x \\ 
                1   &=& x 
\end{array}
$$
Um die Stelle $x=1$ untersuchen wir das Monotonieverhalten der Funktion mittels Vorzeichentabelle:
$$
\renewcommand{\arraystretch}{1.5}
\begin{array}{r|c|c|c}
        & x<1 & x=1 & x>1 \\
\hline 
e^{1-x} &  +  &  +  &  +  \\    
(1-x)   &  +  &  0  &  -  \\
\hline 
f'(x)   &  +  &  0  &  -  \\
\hline 
f(x)    & \nearrow & \rightarrow & \searrow \\
        &  \textnormal{sms} & H  &  \textnormal{smf}     
\end{array}
$$

\noindent Die Funktion ist streng monoton steigend für $x \in ]-\infty ; 1[$. \\
Die Funktion ist streng monoton fallend für $x \in ]1; \infty[$.\\
Der Hochpunkt an der Stelle $x=1$ hat die y-Koordinate:
$$y=f(1)=1 \quad \Rightarrow \quad H(1|1)$$ 

\noindent {\bf Aufgabe 3. c)}\\
Wir leiten die vorgeschlagene Stammfunktion $F$ ab und vergleichen sie mit $f$:
$$
\renewcommand{\arraystretch}{1.5}
\begin{array}{rcl}
F'(x)&=& e^{1-x} - (1-x)\cdot e^{1-x}\\
     &=& (1-(1-x))\cdot e^{1-x} \\
     &=& x \cdot e^{1-x} = f(x)
\end{array}
$$
Die Nullstelle von $F$ ist bei:
$$
\renewcommand{\arraystretch}{1.5}
\begin{array}{rcl}
-e^{1-x}- x \cdot e^{1-x} &=& 0 \\
           e^{1-x} (-1-x) &=& 0 \qquad| : e^{1-x}\\
                   (-1-x) &=& 0 \\
                      x   &=& -1
\end{array}
$$

\begin{center}
	\includegraphics[width=7 cm]{../../viecher/endcomic.pdf}
	
	Hier geht es zurück zum \href{https://www.okuyakl.de/math/m11kdL062/aa062.pdf}{Aufgabenblatt}
\end{center}


\end{document}

